% Part: normal-modal-logic
% Chapter: frame-correspondence
% Section: equivalence-S5

\documentclass[../../../include/open-logic-section]{subfiles}

\begin{document}

\olfileid[fr]{nml}{frd}{es5}

\olsection{Relations d'équivalence et \Log{S5}}

La logique modale \Log{S5} se caractérise comme l'ensemble des
!!{formula}s valides dans tous les cadres universels : chaque monde y
est accessible depuis tout monde, y compris lui-même. Dans ce cas,
$\Box$ correspond à la nécessité et $\Diamond$ à la possibilité :
$\Box !A$ est vraie si $!A$ est vraie dans \emph{tous} les mondes, et
$\Diamond !A$ est vraie si $!A$ est vraie dans \emph{au moins un} monde.
Nous verrons que \Log{S5} se caractérise aussi comme l'ensemble des
!!{formula}s valides dans tous les cadres réflexifs, symétriques et
transitifs, c'est-à-dire ceux dont la relation d'accessibilité est
une \emph{relation d'équivalence}.

\begin{defn}
  Une relation binaire~$R$ sur~$W$ est une \emph{relation
  d'équivalence} si et seulement si elle est réflexive, symétrique et
  transitive. Une relation~$R$ sur~$W$ est \emph{universelle} si et
  seulement si $Ruv$ pour tous $u,v \in W$.
\end{defn}

Puisque \Ax{T}, \Ax{B} et \Ax{4} caractérisent respectivement les
cadres réflexifs, symétriques et transitifs, les cadres dont la
relation d'accessibilité est une relation d'équivalence sont
exactement ceux où ces trois !!{formula}s sont valides. Les relations
d'équivalence se caractérisent aussi par d'autres combinaisons de
!!{formula}s : les conditions de leur définition sont équivalentes à
d'autres combinaisons de propriétés connues de~$R$.

\begin{prop}\ollabel{prop:equivalences}
  Les conditions suivantes sont équivalentes :
  \begin{enumerate}
  \item $R$ est une relation d'équivalence;
  \item $R$ est réflexive et euclidienne;
  \item $R$ est sérielle, symétrique et euclidienne;
  \item $R$ est sérielle, symétrique et transitive.
  \end{enumerate}
\end{prop}

\begin{proof}
  Exercice.
\end{proof}

\begin{prob}
  Démontrer la \olref[nml][frd][es5]{prop:equivalences} en établissant que :
  \begin{enumerate}
  \item si $R$ est symétrique et transitive, elle est euclidienne;
  \item si $R$ est réflexive, elle est sérielle;
  \item si $R$ est réflexive et euclidienne, elle est symétrique;
  \item si $R$ est symétrique et euclidienne, elle est transitive;
  \item si $R$ est sérielle, symétrique et transitive, elle est réflexive.
  \end{enumerate}
  Expliquer pourquoi ces résultats suffisent à démontrer l'équivalence
  des conditions.
\end{prob}

La \olref{prop:equivalences} est le pendant sémantique de la
\olref[prf][prs]{prop:S5} : il donne une autre caractérisation de la
logique modale des cadres sur lesquels $R$ est une relation
d'équivalence, traditionnellement appelée \Log{S5}.

Quel rapport y a-t-il entre relations universelles et relations
d'équivalence? Toute relation universelle est une relation
d'équivalence, mais la réciproque est fausse. Pourtant, les !!{formula}s
valides dans tous les cadres à relation universelle sont exactement
celles qui sont valides dans tous les cadres à relation d'équivalence.

\begin{prop}
  Soit $R$ une relation d'équivalence. Pour chaque $w \in W$,
  définissons la \emph{classe d'équivalence} de $w$ par
  $[w] = \{w'\in W : Rww'\}$. Alors :
  \begin{enumerate}
  \item $w \in [w]$;
  \item $R$ est universelle sur chaque classe d'équivalence~$[w]$;
  \item les classes d'équivalence forment une partition de~$W$ en
    sous-ensembles deux à deux disjoints dont la réunion est~$W$.
  \end{enumerate}
\end{prop}

\begin{prop}\ollabel{prop:S5=univ}
  !!^a{formula}~$!A$ est valide dans tous les cadres
  $\mModel{F} =\tuple{W,R}$ où $R$ est une relation d'équivalence si et
  seulement si elle est valide dans tous les cadres
  $\mModel{F} =\tuple{W,R}$ où $R$ est universelle. La logique des
  cadres universels est donc précisément \Log{S5}.
\end{prop}

\begin{proof}
  Une relation universelle~$R$ sur~$W$ est évidemment une relation
  d'équivalence. Si $!A$ est valide dans tous les cadres où $R$ est
  une relation d'équivalence, elle est donc valide dans tous les
  cadres universels. Pour l'autre sens, raisonnons par contraposée.
  Supposons qu'une !!{formula}~$!B$ soit fausse en un monde~$w$ d'un
  modèle $\mModel{M} = \tuple{W, R, V}$ fondé sur un cadre
  $\tuple{W,R}$ où $R$ est une relation d'équivalence sur~$W$.
  Ainsi, $\mSat/{M}{!B}[w]$. Définissons le modèle
  $\mModel{M}' = \tuple{W', R', V'}$ de la manière suivante :
  \begin{enumerate}
  \item $W' = [w]$;
  \item $R'$ est universelle sur $W'$;
  \item $V'(p) = V(p) \cap W'$. 
  \end{enumerate}
  (L'ensemble~$W'$ des mondes de $\mModel{M}'$ correspond à la région
  grisée de la \olref{fig:partition}.) Les relations $R$ et $R'$
  coïncident sur~$W'$. Une récurrence sur la construction des
  !!{formula}s montre alors que, pour tout $w' \in W'$ et toute
  formule~$!A$, $\mSat{M'}{!A}[w']$ si et seulement si
  $\mSat{M}{!A}[w']$; les deux expressions ont un sens puisque
  $W' \subseteq W$. En particulier, $\mSat/{M'}{!B}[w]$ : la formule
  $!B$ est donc fausse dans un modèle fondé sur un cadre universel.
\end{proof}

\begin{figure}[t]
  \centering
  \begin{tikzpicture}[node distance=2cm, auto, thick]
    \clip [rounded corners] (0,0) -- (6,0) -- (6,4) -- (0,4) --  cycle;
    \begin{scope}
      \clip (0,0) -- (2,0)  .. controls (1.5,1.5) and (4.5,1.5)
      .. (4,4) -- (0,4) -- (0,0) -- cycle;
      \filldraw[fill=gray!40] (0,4) -- (0,2) .. controls (1.5,1.5)
      and (4.5,1.5) .. (4.5,0) -- (6,0) -- (6,4) -- (0,4) --  cycle;
    \end{scope}
    \draw [name path=line1] (2,0) .. controls (1.5,1.5) and (4.5,1.5) .. (4,4);
    \draw [name path=line2] (0,2)  .. controls (1.5,1.5) and (4.5,1.5) .. (4.5,0);
    \path [name intersections={of=line1 and line2}];
    \draw [rounded corners, use as bounding box, very thick] (0,0)
    -- (6,0) -- (6,4) -- (0,4) --  cycle; 
    \node at (2,3) {$[w]$} ;
    \node at (1,1) {$[u]$} ;
    \node at (3,0.75) {$[v]$} ;
    \node at (4.75,2) {$[z]$} ;
  \end{tikzpicture}
  \caption{Une partition de $W$ en classes d'équivalence.}
  \ollabel{fig:partition}
\end{figure}

\end{document}

